\section{Task, Reference, and Relational Evidence}
\label{sec:method}
\paragraph{Source graph and historical scope.}
We reuse the corporate KG snapshot described by Kirpichev et al.~\cite{kirpichev2026towards}, covering S\&P~500 companies and associated sectors, industries, executives, investors, products, resources, locations, and news. The authors confirm unchanged snapshot reuse. We analyze the archived, frozen tasks rather than reconstruct their historical database extraction from current software. Appendix~\ref{app:cypher} reproduces the author-supplied Cypher extraction specification. It defines fixed small instances and an induced-edge filter; it is not presented as the complete executable pipeline for all archived sizes or per-cell sampling quotas. No database query was executed during this analysis.

\subsection{Directed Task Families}
\label{sec:motifs}
We study seven directed families: chain, star-in, star-out, tree-in, tree-out, rhombus, and hourglass. Figure~\ref{fig:raw_motifs} gives a compact key alongside Raw/Raw+E recall. Arrows denote factual KG relations; incoming queries inspect them in the reverse direction. Each query names an anchor, which belongs to the all-route selection reference together with the endpoints and intermediate entities. Stars use one-link incoming or outgoing queries; trees and rhombi use two-link routes, ending earlier only at a dead end. Chains stop at the first requested entity type, and hourglass queries request incoming and outgoing one-link groups separately. The Core sizes are six nodes for chain, stars, and trees, four for rhombus, and five for hourglass. Full graph definitions, traversal rules, and the distinction between endpoint answers and selected route entities appear in Appendix~\ref{app:motifs}. The families also differ in query operations and entity/relation types.


\paragraph{Candidates and distractors.}
Each main task presents 25 named candidate entities, including the motif entities and distractors. Disconnected distractors have no edges to reference entities in either direction, although they may connect to each other. In the adjacent variant, up to ten distractors are neighbours of motif entities; remaining distractors are disconnected. Paired noise variants retain the question and reference but may contain different distractor identities. Explicit evidence includes relations involving distractors, not only reference-induced edges.

\paragraph{Selection target.}
Let $C$ denote the candidate catalogue and $G_C$ its graph-defined reference context, which need not be recoverable from native cards. Let $q$ be the endpoint query, $a$ its anchor, and $\mathcal{P}(q,G_C)$ its allowed supporting routes under the saved traversal/stopping rule. We evaluate the entity-route reference
\begin{equation}
 W(q,G_C)=\{a\}\ \cup\!\bigcup_{p\in\mathcal{P}(q,G_C)} V(p),
\end{equation}
where $V(p)$ contains every entity on route $p$. Endpoint answers $A(q,G_C)$ are distinct from $W$. Simple routes do not revisit an entity; the saved task specification fixes the allowed direction, depth, and terminal-type rule. In the audited structured archive, saved motif/reference IDs coincide with this witness union. Equality is a property of these accepted tasks, not an assumed identity for arbitrary candidate expansion.

For example, one archived rhombus has two branches from the anchor through two intermediate entities to one endpoint. Its endpoint answer has one entity and its reference four. Removing the endpoint card's facts while retaining the name/type catalogue leaves that endpoint recoverable from incoming facts. Thus reference membership does not establish that every card contributes indispensable information. The selection objective is complete route coverage, not minimum sufficient evidence or final answer correctness.

\paragraph{Evidence conditions.}
Raw presents original entity metadata; Raw+E adds explicit relations to the same cards. Actual archived inputs, questions, candidate order, and references match after removing only the relation fields. This is an information-availability contrast; native descriptions are not assumed to express all graph relations. Controlled triples and templated prose instead encode designated matched graph facts. Source-card, destination-card, and both-card variants alter presentation location while retaining relation direction. Names/types remain visible in all conditions. Pseudonyms, direct questions, and linked-attribute encodings provide further diagnostic contrasts detailed in Appendix~\ref{app:conditions}.
